
\section{相关游戏与竞品简要分析}

\textbf{吸血鬼幸存者（Vampire Survivors）}：极简操作 + 自动攻击 + 海量敌人，证明 Survivor-like 品类的市场验证\cite{gregory2018game}。本项目借鉴其战斗节奏与升级构筑，但通过 ECS 与 JSON 配表实现更透明的逻辑扩展，便于在论文中说明模块边界与数据流。

\textbf{杀戮尖塔（Slay the Spire）}：卡牌 Roguelike 与节点地图结合，强调路线决策与风险收益\cite{gregory2018game}。本项目跑图采用类似的 DAG 结构，但战斗为实时俯视角而非回合制，技术难点落在帧率与碰撞而非回合状态机。

\textbf{Noita}：法术模块组合带来高自由度构筑\cite{gregory2018game}。本项目 effect 链式修饰器（修饰器须位于子弹效果之前）吸收了“组合深度”思想，但未实现完整物理模拟与像素级交互，以控制毕业设计范围。

\textbf{其他 H5 生存类作品}：多数采用 Cocos/Laya 预制体配合传统面向对象层次结构。本项目以轻量 ECS、配表与 Worker 作为差异化工程路径，侧重架构可维护性与性能可测性，而非与商业产品在内容量上对标。

\section{术语与缩写预备说明}

全文首次出现英文缩写时给出中文释义，详见附录~\ref{app:terms}。引擎 API 与核心类名保留英文，便于与源码对照。
